\documentclass[a4paper,11pt]{article}
\usepackage[utf8]{inputenc}
\usepackage[margin=1in]{geometry}
\usepackage{amsmath}
\usepackage{graphicx}
\usepackage{listings}
\usepackage{xcolor}
\usepackage{hyperref}
\hypersetup{
	colorlinks=true,
	linkcolor=blue,
	urlcolor=blue,
	citecolor=blue
}

\definecolor{codegray}{rgb}{0.5,0.5,0.5}
\definecolor{backcolour}{rgb}{0.95,0.95,0.92}

\lstdefinestyle{mystyle}{
	backgroundcolor=\color{backcolour},
	commentstyle=\color{codegray},
	keywordstyle=\color{blue},
	numberstyle=\tiny\color{codegray},
	stringstyle=\color{red},
	basicstyle=\ttfamily\footnotesize,
	breaklines=true,
	captionpos=b,
	keepspaces=true,
	numbers=left,
	numbersep=5pt,
	showspaces=false,
	showstringspaces=false,
	showtabs=false,
	tabsize=2
}

\lstset{style=mystyle}

\title{Brief Introduction to Python}
\author{Tiago Castro}
\date{\today}

\begin{document}
	
	\maketitle
	
	\section{Variables}
	In Python, variables are dynamically typed. A variable can hold values of any type, and its type may change during execution.
	
	\begin{lstlisting}[language=Python]
		x = 5          # Integer
		name = "Alice"  # String
		pi = 3.1415     # Float
		is_valid = True  # Boolean
	\end{lstlisting}
	
	\subsection*{Exercises}
	\begin{enumerate}
		\item Declare a variable `age` with your age and print its type using the `type` function.
		\item Assign a string to a variable `greeting` and concatenate it with a name variable.
		\item Change the type of `x` from integer to string and print both values.
	\end{enumerate}
	
	\section{Lists, Dictionaries and Sets}
	\subsection{Lists}
	A list is an ordered and mutable collection of elements. Elements can be of mixed types.
	
	\begin{lstlisting}[language=Python]
		fruits = ["apple", "banana", "cherry"]
		print(fruits[1])          # Output: banana
		fruits.append("date")    # Add an element
		fruits[0] = "avocado"    # Modify an element
	\end{lstlisting}
	
	\subsection*{Exercises}
	\begin{enumerate}
		\item Create a list of five numbers and compute its sum.
		\item Reverse a list `colors = ["red","green","blue"]` and print the result.
		\item Remove the last element from a list and print the remaining elements.
	\end{enumerate}
	
	\subsection{Dictionaries}
	A dictionary stores data as key-value pairs. Keys must be unique and immutable.
	
	\begin{lstlisting}[language=Python]
		person = {"name": "Alice", "age": 25}
		print(person["name"])     # Output: Alice
		person["city"] = "Rome"  # Add new key-value pair
	\end{lstlisting}
	
	\subsection*{Exercises}
	\begin{enumerate}
		\item Create a dictionary mapping three countries to their capitals.
		\item Update the dictionary by changing one capital and print the changes.
		\item Write a loop to print all keys and values of the dictionary.
	\end{enumerate}
	
	\subsection{Sets}
	A set is an unordered collection of unique elements. Useful for membership tests and eliminating duplicates.
	
	\begin{lstlisting}[language=Python]
		unique_numbers = {1, 2, 3, 2}
		print(unique_numbers)       # Output: {1, 2, 3}
		unique_numbers.add(4)       # Add an element
	\end{lstlisting}
	
	\subsection*{Exercises}
	\begin{enumerate}
		\item Create two sets and compute their union, intersection and difference.
		\item Remove a specific element from a set.
		\item Check if one set is a subset of another.
	\end{enumerate}
	
	\section{List Comprehension}
	List comprehensions provide a concise syntax to create lists based on existing iterables.
	
	\begin{lstlisting}[language=Python]
		squares = [x**2 for x in range(5)]
		print(squares)             # Output: [0, 1, 4, 9, 16]
		evencubes = [x**3 for x in range(10) if x % 2 == 0]
		print(even_cubes)          # Output: [0, 8, 64, 216, 512]
	\end{lstlisting}
	
	\subsection*{Exercises}
	\begin{enumerate}
		\item Use a list comprehension to extract the first letter of each word in `words = ["hello","world"]`.
		\item Create a list of tuples `(x, x*2)` for `x` in range 5.
		\item Flatten a nested list `[[1,2],[3,4],[5]]` using a single comprehension.
	\end{enumerate}
	
	\section{Functions}
	Functions allow organization of code into reusable blocks. Define a function with `def`, include parameters and return values.
	
	\begin{lstlisting}[language=Python]
		def greet(name):
		"""Return a greeting for the given name."""
		return f"Hello, {name}!"
		
		print(greet("Alice"))  # Output: Hello, Alice!
		
		# Function with default parameter
		
		def power(base, exponent=2):
		"""Return base raised to the given exponent."""
		return base ** exponent
		
		print(power(3))         # Output: 9
		print(power(2, 3))      # Output: 8
		
		# Function with arbitrary arguments
		
		def summarize(*args):
		"""Return sum of all arguments."""
		return sum(args)
		
		print(summarize(1, 2, 3, 4))  # Output: 10
	\end{lstlisting}
	
	\subsection*{Exercises}
	\begin{enumerate}
		\item Write a function `square(x)` that returns the square of x.
		\item Create a function that takes a list and returns its maximum value.
		\item Define a function that accepts keyword arguments and prints them.
	\end{enumerate}
	
	\section{NumPy}
	NumPy is a library for efficient numerical computing with support for large, multi-dimensional arrays and matrices.
	
	\begin{lstlisting}[language=Python]
		import numpy as np
		
		array = np.array([1, 2, 3])
		print(array * 2)            # Output: [2 4 6]
		
		matrix = np.array([[1, 2], [3, 4]])
		print(matrix.T)             # Transpose
		print(np.dot(matrix, matrix))  # Matrix multiplication
	\end{lstlisting}
	
	\subsection*{Exercises}
	\begin{enumerate}
		\item Create a 3x3 identity matrix using NumPy.
		\item Compute the mean and standard deviation of `array = np.arange(10)`.
		\item Generate an array of 20 linearly spaced points between 0 and 1.
	\end{enumerate}
	
	\section{Matplotlib}
	Matplotlib is a plotting library for Python. It provides an object-oriented API and state-based interface.
	
	\begin{lstlisting}[language=Python]
		import matplotlib.pyplot as plt
		import numpy as np
		
		x = np.linspace(0, 2* np.pi, 100)
		y = np.sin(x)
		
		plt.figure()
		plt.plot(x, y)
		plt.xlabel('x')
		plt.ylabel('sin(x)')
		plt.title('Sine Wave')
		plt.grid(True)
		plt.show()
	\end{lstlisting}
	
	\subsection*{Exercises}
	\begin{enumerate}
		\item Plot a cosine function on the same axes as the sine function.
		\item Create a scatter plot for random points using `np.random.rand(50,2)`.  
		\item Save a plot to file `output.png` without displaying it interactively.
	\end{enumerate}
	
\end{document}
